\chapter{ОЦЕНКА ПРОЕКТА КАК СТАРТАПА: БИЗНЕС-МОДЕЛЬ, РЫНОК И ИНВЕСТИЦИОННАЯ ПРИВЛЕКАТЕЛЬНОСТЬ}

Разработанное в рамках настоящей выпускной квалификационной работы веб-приложение «Shasl MP» рассматривается не только как технический результат, но и как самостоятельный стартап-проект, ориентированный на практическое внедрение и коммерческое использование. В данной главе проект анализируется через призму конкурса «Диплом как стартап»: формулируется проблема и предлагаемое решение, обосновывается уникальность продукта, описывается бизнес-модель, проводится анализ рынка и конкурентов, приводятся финансовые показатели и план роста, оцениваются риски и определяется необходимый объём инвестиций. Цель главы --- показать, что результат работы обладает признаками готового цифрового продукта и имеет потенциал для устойчивого коммерческого развития.

\section{Описание проблемы и предлагаемого решения}

Электронная коммерция в России демонстрирует устойчивый рост, а количество продавцов (селлеров) на маркетплейсах Ozon и Wildberries ежегодно увеличивается. Вместе с тем по мере роста оборота продавцы сталкиваются с возрастающей сложностью управления товарными запасами и аналитикой продаж. Маркетплейсы предоставляют отчёты в форматах CSV и XLSX, однако эти данные слабо структурированы и не содержат готовых управленческих рекомендаций.

Ключевая проблема, выявленная в ходе работы, заключается в следующем: продавцы на маркетплейсах тратят до 70 \% рабочего времени на ручную обработку отчётов и не имеют инструмента, который автоматически прогнозирует, когда и в каком регионе закончится товар, а также формирует рекомендации по пополнению остатков. Это приводит к ряду негативных последствий:

\begin{itemize}
\item высокая вероятность ошибок при ручной обработке данных в Excel;
\item финансовые издержки вследствие неполного или неточного анализа;
\item невозможность оперативно реагировать на изменение остатков на складах;
\item значительные затраты времени, которое могло бы быть направлено на развитие бизнеса.
\end{itemize}

Предлагаемое решение --- веб-приложение «Shasl MP», которое автоматизирует аналитику продаж и управление товарными остатками. Система реализует три взаимосвязанных направления функциональности:

\begin{enumerate}
\item автоматизированное получение данных через API маркетплейсов Ozon и Wildberries без необходимости ручной выгрузки отчётов;
\item расчёт ключевых метрик и показателей бизнеса --- юнит-экономики, дней до исчерпания запаса, кластерного анализа остатков, динамики продаж, выручки, индекса локализации;
\item автоматическую структуризацию данных логикой программы и предоставление наглядной визуализации в веб-интерфейсе.
\end{enumerate}

Таким образом, продукт переводит продавца от трудоёмкой ручной аналитики к принятию обоснованных управленческих решений на основе актуальных данных в реальном времени. Девиз проекта --- «Доступная аналитика каждому».

\section{Уникальность продукта}

Уникальность проекта основана на смещении фокуса аналитики с вопроса «сколько товара осталось» на вопрос «когда и где товар закончится». Большинство существующих сервисов отображают лишь общий остаток, тогда как «Shasl MP» строит адаптивный кластерный прогноз по каждому складу и региону, формируя готовые рекомендации по пополнению --- «куда и сколько везти».

Для проверки ключевой гипотезы --- что продавцы нуждаются не просто в отчётах, а в готовых рекомендациях по пополнению остатков --- было проведено восемь глубинных интервью с действующими селлерами Ozon и Wildberries с объёмом продаж от 100 тыс. до 15 млн руб./мес. По результатам интервью выявлены следующие основные «боли» целевой аудитории:

\begin{itemize}
\item 100 \% опрошенных тратят от 4 до 10 часов в неделю на ручную аналитику в Excel;
\item 75 \% сталкивались с упущенной выручкой из-за того, что товар закончился в одном регионе, но оставался невостребованным («мёртвым грузом») в другом;
\item 87 \% не используют готовые аналитические сервисы из-за высокой стоимости и отсутствия нужных метрик, прежде всего прогноза по кластерам.
\end{itemize}
Для проверки решения был разработан минимально жизнеспособный продукт (MVP) с ключевым модулем --- кластерным прогнозом остатков. MVP протестирован на трёх продавцах. Полученная обратная связь подтвердила ценность продукта: пользователи отметили, что ранее не знали об исчерпании товара в конкретном регионе, пока не получали негативные отзывы, тогда как функция «дней до исчерпания» по каждому складу позволяет видеть ситуацию заблаговременно --- за 30 дней. По итогам тестирования ключевая гипотеза подтверждена, а уникальность продукта определена как переход от констатации остатка к прогнозированию его исчерпания в разрезе складов.

Дополнительным фактором уникальности является технологическая реализация. В отличие от ранней версии приложения, требовавшей ручной загрузки CSV-отчётов, в рамках работы реализована полноценная интеграция с API маркетплейсов, выбор временных периодов непосредственно в интерфейсе, а также безопасное хранение и использование API-ключей. Это позволяет автоматически сравнивать поставки по схемам FBO и FBS за произвольный период и формировать аналитику в реальном времени, что существенно снижает порог входа для продавцов, не обладающих специальными навыками работы с данными.

\section{Результаты тестирования и подтверждение спроса}

Перед коммерческим запуском продукт прошёл практическую проверку на реальных пользователях. Минимально жизнеспособный продукт включал подсистему авторизации, API-интеграцию с Ozon, кластерный прогноз остатков с расчётом дней до исчерпания по каждому складу и визуализацию продаж. Методика сбора данных сочетала качественные и количественные методы: глубинные интервью с восемью селлерами, логирование действий пользователей в приложении, замеры времени аналитики до и после внедрения, а также форму обратной связи, встроенную в интерфейс.

Тестирование проводилось в течение трёх месяцев с поэтапным расширением числа участников. Динамика результатов по месяцам приведена в таблице~\ref{tab:testResults}.

\begin{table}[H]
\caption{Результаты тестирования MVP по месяцам}
\label{tab:testResults}
\small
\begin{center}
\vspace{-0.5cm}\begin{tabular}{|p{1.5cm}|p{3cm}|p{8cm}|}
\hline
Период & Тестировщиков & Ключевые результаты \\ \hline
1-й месяц & 3 & API-интеграция стабильна; выявлена и устранена проблема кодировки Windows-1251 \\ \hline
2-й месяц & 5 & Кластерный прогноз признан главной ценностью; время аналитики сократилось с 1-3 часов до 15-20 минут в неделю \\ \hline
3-й месяц & 7 & 100 процентов участников готовы рекомендовать продукт; 6 из 7 пользуются приложением ежедневно \\ \hline
\end{tabular}
\end{center}
\end{table}

Полученные результаты подтверждают как работоспособность продукта, так и наличие устойчивого спроса. Ключевым измеримым эффектом стало многократное сокращение трудозатрат на аналитику --- с нескольких часов до 15-20 минут в неделю, что напрямую отражает создаваемую продуктом ценность. Высокая доля пользователей, готовых рекомендовать продукт, и ежедневный характер использования свидетельствуют о хорошей удерживающей способности (retention) и формируют основу для органического роста за счёт «сарафанного радио» и реферальной программы.

\section{Бизнес-модель}

Бизнес-модель проекта построена на модели подписки (SaaS) и описана с использованием структуры бизнес-модели Остервальдера. Основные элементы модели представлены в таблице~\ref{tab:businessModel}.

\begin{table}[H]
\caption{Бизнес-модель проекта «Shasl MP»}
\label{tab:businessModel}
\begin{center}
\begin{tabular}{|l|p{10cm}|}
\hline
Компонент & Описание \\ \hline
Ценностное предложение & Автоматическая выгрузка отчётов через API, визуализация продаж по дням, кластерный прогноз остатков с рекомендациями (FBO/FBS), автоматический расчёт прибыли с подключением юнит-экономики \\ \hline
Потоки доходов & Подписка: тариф «Базовый» (адаптивный блок с продажами, аналитикой и юнит-экономикой); тариф «Продвинутый» (всё перечисленное, а также кластерный анализ, мультиаккаунт, встроенный ИИ-агент для аналитики) \\ \hline
Ключевые ресурсы & Сервер, API-интеграция, алгоритмы веб-приложения \\ \hline
Сегменты клиентов & Селлеры с оборотом от 100 тыс. руб./мес. \\ \hline
Каналы сбыта & Сайт с демо-доступом, Telegram-чат, партнёры (агрегаторы), медийная реклама с демонстрацией использования \\ \hline
Издержки & Хостинг, API-запросы, разработка, поддержка, заработная плата сотрудников \\ \hline
Масштабирование & Подключение новых маркетплейсов (WB, Яндекс Маркет), модуль автозаказа \\ \hline
\end{tabular}
\end{center}
\end{table}

Монетизация осуществляется по подписочной модели со средней стоимостью около 500 руб./мес. на пользователя. Двухуровневая тарифная сетка позволяет охватить как начинающих продавцов, так и более крупных селлеров, заинтересованных в расширенной аналитике, мультиаккаунте и интеллектуальных функциях. Каналы привлечения сочетают органический трафик (демо-доступ на сайте, тематическое сообщество в Telegram) и платные каналы (медийная реклама, партнёрские размещения у агрегаторов).

Целевая аудитория продукта неоднородна и условно делится на три группы. Первая --- начинающие продавцы с оборотом от 100 тыс. руб./мес., которым важна простота интерфейса и базовая аналитика по доступной цене; для них предназначен тариф «Базовый». Вторая --- растущие продавцы со средним оборотом, для которых критичен кластерный прогноз остатков и управление поставками по схемам FBO/FBS; они формируют ядро платящей аудитории тарифа «Продвинутый». Третья --- крупные селлеры и небольшие агентства, управляющие несколькими аккаунтами, для которых ценны мультиаккаунт и ИИ-агент аналитики. Такая сегментация позволяет выстраивать дифференцированное ценообразование и последовательно повышать пожизненную ценность клиента по мере роста его бизнеса.

Стратегия выхода на рынок предусматривает запуск с бесплатным демо-доступом, снимающим барьер недоверия к новому сервису, с последующей конверсией в платную подписку после демонстрации ценности кластерного прогноза. Сообщество в Telegram выполняет роль канала поддержки и площадки для сбора обратной связи, а партнёрские размещения у агрегаторов и инфлюенсеров в нише электронной коммерции обеспечивают приток целевой аудитории. Реферальная программа стимулирует действующих пользователей привлекать новых, снижая совокупную стоимость привлечения клиента.

\section{Анализ рынка}

Целевой рынок проекта --- продавцы российских маркетплейсов Ozon и Wildberries, испытывающие потребность в автоматизированной аналитике продаж и управлении остатками. Рынок характеризуется высокими темпоми роста электронной коммерции и ежегодным увеличением числа активных селлеров, что напрямую расширяет потенциальную клиентскую базу сервисов аналитики.

Подтверждением спроса служат результаты проведённого исследования целевой аудитории. Проверка ключевой гипотезы на выборке из восьми действующих селлеров показала, что потребность в готовых рекомендациях по пополнению остатков и адаптивной визуализации данных носит массовый характер: все опрошенные тратят значительное время на ручную аналитику, а большинство сталкивались с прямыми финансовыми потерями из-за несбалансированного распределения товара между регионами. При этом существующие сервисы не покрывают потребность в кластерном прогнозе, что формирует свободную нишу для предлагаемого продукта.

Сегментация рынка проводится по обороту продавца. Базовый целевой сегмент --- селлеры с оборотом от 100 тыс. руб./мес., для которых стоимость подписки является несущественной по сравнению с потенциальной экономией времени и предотвращёнными потерями выручки. По мере развития продукта рынок расширяется за счёт подключения новых торговых площадок (Яндекс Маркет) и выхода на рынки СНГ (Казахстан, Беларусь), что многократно увеличивает совокупный адресуемый рынок.

Оценку рынка целесообразно проводить по модели TAM-SAM-SOM. Совокупный адресуемый рынок (TAM) формируют все продавцы российских маркетплейсов, число которых исчисляется сотнями тысяч и продолжает расти. Доступный рынок (SAM) ограничивается селлерами с оборотом от 100 тыс. руб./мес., для которых аналитика является значимой статьёй принятия решений. Реально достижимый рынок (SOM) на горизонте трёх лет оценивается в десятки тысяч платящих пользователей, что соответствует целевым показателям плана роста. Дополнительным драйвером спроса выступает усиление конкуренции между самими продавцами: по мере насыщения ниш выигрывают те селлеры, кто принимает решения на основе данных, а не интуиции, что повышает ценность аналитических инструментов.

\section{Анализ конкурентов}

На рынке аналитических сервисов для маркетплейсов присутствует ряд игроков, среди которых наиболее значимыми конкурентами являются MPStats, Moneyplace и Stat4market. Эти сервисы предоставляют широкий набор аналитических инструментов, однако обладают существенными недостатками с точки зрения целевой аудитории проекта: высокой стоимостью и отсутствием специализированного кластерного прогноза остатков по складам.

Конкурентные преимущества «Shasl MP» состоят в следующем:

\begin{itemize}
\item визуализация продаж по дням --- графики с детализацией по дням и по схемам поставок (FBO, FBS);
\item анализ динамики продаж в разрезе каждой схемы поставок;
\item кластерный анализ --- адаптивный прогноз и визуализация остатков и спроса товаров по регионам;
\item автоматический расчёт прибыли с подключением модуля юнит-экономики.
\end{itemize}

Сводное сопоставление продукта с конкурентами по ключевым параметрам приведено в таблице~\ref{tab:competitors}.

\begin{table}[H]
\caption{Сравнение с конкурентами}
\label{tab:competitors}
\begin{center}
\vspace{-0.5cm}\begin{tabular}{|p{4cm}|p{3cm}|p{3cm}|p{3cm}|}
\hline
Параметр & Shasl MP & Крупные сервисы & Excel вручную \\ \hline
Кластерный прогноз остатков & Да & Нет & Нет \\ \hline
Рекомендации по пополнению & Да & Частично & Нет \\ \hline
Автоматизация через API & Да & Да & Нет \\ \hline
Юнит-экономика & Да & Частично & Частично \\ \hline
Стоимость для селлера & Низкая & Высокая & Нулевая \\ \hline
Затраты времени & Минимальные & Средние & Высокие \\ \hline
\end{tabular}
\end{center}
\end{table}

Для закрепления лидерства запланированы следующие шаги по развитию продукта: расширение функционала анализа продаж за счёт новых аналитических возможностей; подключение блока управления рекламой и её анализа; внедрение ИИ-агента для продвинутого анализа; расширенная работа с карточкой товара --- поиск функциональных ключей и запросов, написание SEO-описаний, аналитика карточки, анализ конкурентов и выбор ниш. Эти направления формируют долгосрочную дорожную карту и затрудняют копирование продукта конкурентами.

\section{Финансовые показатели}

Финансовая модель проекта основана на подписке со средней стоимостью около 500 руб./мес. на одного пользователя. Юнит-экономика рассчитана в расчёте на одного платящего пользователя, что позволяет оценить маржинальность и точку безубыточности. Ключевые финансовые параметры проекта приведены в таблице~\ref{tab:finance}.

\begin{table}[H]
\caption{Ключевые финансовые показатели}
\label{tab:finance}
\begin{center}
\begin{tabular}{|l|c|}
\hline
Показатель & Значение \\ \hline
Средняя стоимость подписки & приблизительно 500 руб./мес. \\ \hline
Точка безубыточности & приблизительно 250 пользователей (5-6-й месяц) \\ \hline
Объём инвестиций & 1 000 000 руб. за 15 \% доли компании \\ \hline
Срок окупаемости & 12-14 месяцев \\ \hline
\end{tabular}
\end{center}
\end{table}

Доходная часть модели формируется за счёт ежемесячных подписных платежей и масштабируется пропорционально числу платящих пользователей. Расходная часть включает постоянные и переменные издержки: хостинг и серверную инфраструктуру, оплату API-запросов, разработку и поддержку продукта, а также заработную плату сотрудников. По мере роста числа пользователей доля переменных издержек на одного пользователя снижается, что обеспечивает рост маржинальности при масштабировании.

Расчёты показывают, что при достижении уровня примерно 250 платящих пользователей проект выходит на точку безубыточности, ожидаемую на 5-6-й месяц после запуска коммерческой версии. Полная окупаемость вложенных инвестиций прогнозируется в горизонте 12-14 месяцев. Таким образом, проведённый анализ юнит-экономики подтверждает потенциальную экономическую эффективность проекта и возможность его коммерческого внедрения.

Важными показателями подписочной модели являются пожизненная ценность клиента (LTV) и стоимость его привлечения (CAC). По мере развития продукта планируется увеличение LTV до 8000 руб. за счёт расширения функциональности и повышения удержания, тогда как стоимость привлечения снижается благодаря органическим каналам --- демо-доступу, рекомендациям пользователей и реферальной программе. Положительное соотношение LTV/CAC является ключевым условием устойчивой экономики проекта и инвестиционной привлекательности. Высокая доля постоянных издержек на ранней стадии при низких переменных затратах на одного пользователя обеспечивает выраженный эффект масштаба: каждый дополнительный платящий пользователь после прохождения точки безубыточности приносит преимущественно маржинальный доход.

\section{План роста и масштабирования}

План развития проекта рассчитан на трёхлетний горизонт и предусматривает поэтапный рост клиентской базы, выручки и капитализации. Каждый этап сопровождается достижением конкретных продуктовых и финансовых целей. Дорожная карта развития представлена в таблице~\ref{tab:growthPlan}.

\begin{table}[H]
\caption{План роста проекта на 2025-2028 гг.}
\label{tab:growthPlan}
\begin{center}
\small
\vspace{-0.5cm}\begin{tabular}{|p{3cm}|p{4cm}|p{4cm}|p{4cm}|}
\hline
Период & Результаты / выручка & Ключевые цели & Следующие шаги \\ \hline
1 год (2025-2026) & MVP запущен, приложение опубликовано, интеграция с Ozon и Wildberries; приблизительно 2000 платных пользователей; выручка приблизительно 1 000 000 руб. & Достичь точки безубыточности, удержать 70 \% платящих пользователей & Масштабирование маркетинга, запуск реферальной программы, сбор фидбэка, доработка кластерного прогноза \\ \hline
2 год (2026-2027) & приблизительно 10 000 платных пользователей; выручка приблизительно 7 000 000 руб.; положительный денежный поток & Расширить продукт (Яндекс Маркет, WB), увеличить LTV до 8000 руб. & Интеграция с новыми маркетплейсами, внедрение ИИ-агента, запуск модуля автозаказа \\ \hline
3 год (2027-2028) & приблизительно 50 000 платных пользователей; выручка приблизительно 50 000 000 руб. & Выход на рынки СНГ (Казахстан, Беларусь), рост капитализации в 10 раз & Локализация под новые маркетплейсы, международный маркетинг, поиск стратегических партнёров, патентование алгоритма \\ \hline
\end{tabular}
\end{center}
\end{table}

Стратегия роста сочетает экстенсивное расширение (новые маркетплейсы и географические рынки) и интенсивное развитие (увеличение пожизненной ценности клиента --- LTV --- за счёт новых функциональных модулей). На первом этапе приоритет отдаётся достижению безубыточности и удержанию платящих пользователей; на втором --- расширению продукта и выходу на положительный денежный поток; на третьем --- масштабированию на рынки СНГ и кратному росту капитализации. Внедрение реферальной программы и партнёрских каналов обеспечивает снижение стоимости привлечения клиента по мере роста.

\section{Анализ рисков}

Реализация проекта сопряжена с рядом рисков технического, кадрового, рыночного и юридического характера. Для каждого риска определены вероятность наступления, степень влияния на проект и меры контроля и минимизации. Матрица рисков представлена в таблице~\ref{tab:risks}.

\begin{table}[H]
\caption{Анализ рисков проекта}
\label{tab:risks}
\begin{center}
\vspace{-0.5cm}\begin{tabular}{|p{3cm}|p{3cm}|p{3cm}|p{5cm}|}
\hline
Риск & Вероятность & Влияние & Меры минимизации \\ \hline
Изменение API маркетплейсов & Средняя & Высокое & Мониторинг документации, отслеживание уведомлений об изменениях \\ \hline
Уход ключевого разработчика & Низкая & Высокое & Документирование кода, распределение знаний в команде, найм второго разработчика \\ \hline
Утечка данных или компрометация API-ключей & Низкая & Критическое & JWT-токены, HTTP-only cookies, хеширование паролей, переменные окружения, ротация ключей \\ \hline
Снижение платёжеспособности селлеров & Средняя & Среднее & Гибкая ценовая политика, self-hosted-версия, акции для малого бизнеса \\ \hline
Блокировка маркетплейсом & Низкая & Высокое & Соблюдение лимитов запросов, использование официальных API, юридическое сопровождение \\ \hline
\end{tabular}
\end{center}
\end{table}

Наиболее значимым техническим риском является изменение API маркетплейсов, способное нарушить работу интеграции; он минимизируется регулярным мониторингом документации и оперативной адаптацией кода. Критическим по уровню влияния признан риск утечки данных и компрометации API-ключей, для нейтрализации которого реализован комплекс мер информационной безопасности. Рыночные риски, связанные со снижением платёжеспособности продавцов, компенсируются гибкой тарифной политикой и наличием self-hosted-версии. Совокупно матрица рисков демонстрирует, что наиболее вероятные риски имеют управляемый характер, а риски с критическим влиянием обладают низкой вероятностью.

\section{Инвестиции}

Для доведения продукта до полноценной коммерческой версии и привлечения первых платящих пользователей требуется привлечение инвестиций в объёме 1 000 000 руб. в обмен на 15 \% доли в компании. Распределение инвестиций по статьям расходов приведено в таблице~\ref{tab:investments}.

\begin{table}[H]
\caption{Структура использования инвестиций}
\label{tab:investments}
\begin{center}
\begin{tabular}{|l|r|}
\hline
Статья расходов & Сумма, руб. \\ \hline
Доработка MVP до полноценной версии (фронтенд и бэкенд) & 300 000 \\ \hline
Интеграция с Wildberries и Яндекс Маркет & 200 000 \\ \hline
Маркетинг и привлечение первых платных пользователей & 100 000 \\ \hline
Юридическое сопровождение и регистрация ИП/ООО & 100 000 \\ \hline
Резервный фонд & 300 000 \\ \hline
\textbf{Итого} & \textbf{1 000 000} \\ \hline
\end{tabular}
\end{center}
\end{table}

Структура инвестиций сбалансирована между развитием продукта, привлечением клиентов и обеспечением юридической и финансовой устойчивости. Наибольшая доля средств направляется на доработку продукта и расширение интеграций, что напрямую увеличивает ценность сервиса для пользователей. Выделение резервного фонда (30 \% бюджета) обеспечивает финансовую подушку для покрытия непредвиденных расходов на ранней стадии. С учётом прогнозируемого срока окупаемости 12-14 месяцев и потенциала масштабирования предлагаемые условия инвестирования являются привлекательными для венчурного финансирования.

Привлекаемые средства рассматриваются как посевные инвестиции (seed-стадия), необходимые для перехода от проверенного на рынке MVP к полноценному коммерческому продукту и захвата первой доли рынка. Предлагаемая оценка компании в 6 667 тыс. руб. (исходя из 15 \% за 1 000 000 руб.) соответствует ранней стадии проекта с подтверждённым спросом и работающей технологией. Для инвестора привлекательность вложения определяется сочетанием низкого порога входа, короткого срока окупаемости и значительного потенциала роста капитализации: согласно плану развития, к третьему году капитализация проекта может возрасти кратно за счёт роста выручки и выхода на новые рынки. В качестве возможных сценариев выхода (exit) рассматриваются продажа доли стратегическому инвестору из числа крупных игроков рынка e-commerce-аналитики, а также дальнейшие раунды финансирования для ускоренного масштабирования.

\section{Выводы по главе}

В данной главе проект «Shasl MP» был рассмотрен как стартап в формате конкурса «Диплом как стартап». Сформулирована актуальная рыночная проблема и предложено решение в виде специализированного веб-приложения для аналитики и прогнозирования данных маркетплейсов. Уникальность продукта подтверждена результатами глубинных интервью и тестирования MVP и заключается в переходе от констатации остатка к прогнозированию его исчерпания по складам с выдачей готовых рекомендаций.

Разработана бизнес-модель на основе подписки, проведён анализ растущего рынка и конкурентного окружения, в котором продукт занимает свободную нишу кластерного прогноза. Финансовая модель демонстрирует выход на точку безубыточности при 250 пользователях и окупаемость в горизонте 12-14 месяцев. Трёхлетний план роста предусматривает масштабирование до 50 000 платных пользователей и выручки около 50 млн руб. с выходом на рынки СНГ. Выполненный анализ рисков и обоснование инвестиций в объёме 1 000 000 руб. подтверждают, что проект обладает признаками готового цифрового продукта и имеет потенциал для устойчивого коммерческого развития.

Совокупность представленных результатов позволяет заключить, что выпускная квалификационная работа выходит за рамки исключительно учебной задачи и формирует основу полноценного стартап-проекта. Проверенная на реальных пользователях гипотеза, измеримая ценность для клиента, проработанная экономика и масштабируемая технологическая архитектура в совокупности образуют целостное инвестиционное предложение. Дальнейшее развитие проекта связано с расширением перечня поддерживаемых маркетплейсов, внедрением методов машинного обучения для прогнозирования спроса, развитием системы рекомендаций и переходом к облачной модели предоставления сервиса, что соответствует логике конкурса «Диплом как стартап» и подтверждает практическую и коммерческую значимость выполненной работы.

